\documentclass[10pt]{beamer}
\usetheme{metropolis} % Você pode trocar por: Warsaw, Berlin, Metropolis, etc.

\usepackage[utf8]{inputenc}
\usepackage[brazil]{babel}
\usepackage{graphicx}
\usepackage{booktabs}
\usepackage{amsmath}
\usepackage{verbatim}

\title{Detecção Automática de Quedas}
\author{Caio Passos Torkst Ferreira}
\date{\today}

\begin{document}

% Slide de título
\begin{frame}
  \titlepage
\end{frame}

% Sumário
\begin{frame}{Sumário}
  \tableofcontents
\end{frame}

\section{Introdução}

\begin{frame}{Contexto e Motivação}
  \begin{itemize}
    \item Quedas são eventos críticos em ambientes militares.
    \item Detecção automática é vital para resposta rápida.
    \item Machine Learning pode capturar padrões em sensores corporais.
    \item Projeto baseado no artigo:
    \begin{itemize}
      \item \textit{A Machine Learning Approach to Automatic Fall Detection of Soldiers}
      \item \url{https://arxiv.org/abs/2501.15655v2}
    \end{itemize}
  \end{itemize}
\end{frame}



\begin{frame}{Projeto e Dados}
  \begin{itemize}
    \item Reimplementação em PyTorch com melhorias:
    \begin{itemize}
      \item Inclusão de LSTM, K-Fold CV, Optuna (bayesiano), Early Stopping
    \end{itemize}
    \item Pipeline automatizado: \texttt{all.sh}
    \item Dados inerciais do \textbf{peito} (acel. + giro, 3 eixos → shape: (1020, 6))
    \item Total de amostras: \textbf{6038}, Label: \texttt{binary\_two}
    \item Cenário: \texttt{Sc\_4\_T} (tempo)
  \end{itemize}
\end{frame}


\begin{frame}{Coleta de Dados e Sensores}
  \begin{itemize}
    \item 15 voluntários com:
    \begin{itemize}
      \item 1 smartphone (peito), 2 smartwatches (pulsos)
    \end{itemize}
    \item Sensores:
    \begin{itemize}
      \item Aceleração linear (acc) e angular (gyr) em x, y, z
    \end{itemize}
    \item Amostragem média:
    \begin{itemize}
      \item Smartphone: 221 Hz (acc), 219 Hz (gyr)
      \item Smartwatches: 93 Hz (acc/gyr)
    \end{itemize}
  \end{itemize}
\end{frame}


\begin{frame}{Objetivos do Projeto}
  \begin{itemize}
    \item Construir pipeline de detecção de quedas com MLP, CNN1D e LSTM
    \item Validar com K-Fold, Otimização bayesiana (Optuna) e Early Stopping
    \item Avaliar com: \textbf{MCC}, \textbf{F1}, \textbf{Specificity}, \textbf{Accuracy}
    \item Estabilidade: 20 execuções por arquitetura
    \item Interpretabilidade com SHAP e permutação
    \item Estudar impacto do volume de dados nas curvas de aprendizado
  \end{itemize}
\end{frame}

\begin{frame}{Métricas e Justificativa da Abordagem}
  \begin{itemize}
    \item \textbf{F1-Score}: Média harmônica entre precisão e revocação — boa com dados equilibrados.
    \item \textbf{MCC}: Considera TP, TN, FP, FN — ideal para dados desbalanceados (1:8).
    \item Abordagem visa:
    \begin{itemize}
      \item Robustez e confiabilidade em cenários reais
      \item Portabilidade para outros sensores/posições
      \item Explicabilidade em contextos críticos
    \end{itemize}
  \end{itemize}
\end{frame}

\section{Metodologia}

\begin{frame}{Pipeline e Fluxo do Experimento}
  \begin{itemize}
    \item Automação via \texttt{all.sh}; etapas principais:
    \begin{enumerate}
      \item Geração do dataset: \texttt{chest/Sc\_4\_T}
      \item Otimização de hiperparâmetros com \textbf{Optuna} (5-Fold CV, MCC)
      \item Treinamento final com os melhores parâmetros
      \item Avaliação estatística, interpretabilidade (SHAP, Permutation)
    \end{enumerate}
    \item Cada modelo é treinado \textbf{20 vezes} para avaliar estabilidade.
  \end{itemize}
\end{frame}

\begin{frame}{Arquiteturas e Espaços de Busca}
  \begin{itemize}
    \item \textbf{MLP}: 1–3 camadas densas, 32–512 neurônios, dropout [0.0–0.5]
    \item \textbf{CNN1D}: 1–3 convoluções, kernel (3/5/7), filtros (8–64)
    \item \textbf{LSTM}: 1–2 camadas, 16–256 hidden units
    \item \textbf{Learning rate}: [$10^{-4}$, $10^{-2}$] (log)
    \item Saída comum: função \texttt{Sigmoid}
  \end{itemize}
\end{frame}

\begin{frame}{Otimização e Técnicas Aplicadas}
  \begin{itemize}
    \item \textbf{Optuna (TPE)} com 5-Fold CV, métrica alvo: MCC
    \item Técnicas:
    \begin{itemize}
      \item Early Stopping : Evita Overfitting
      \item Median Pruning : Interrompe trials ruins
    \end{itemize}
    \item Parâmetros mais relevantes:
    \begin{itemize}
      \item \textbf{Learning rate} → decisivo em todos os modelos
      \item CNN1D: kernel e filtros
      \item MLP: número de camadas e dropout
      \item LSTM: camadas e hidden size
    \end{itemize}
  \end{itemize}
\end{frame}


\begin{frame}{Treinamento Final}
  \begin{itemize}
    \item Após definição dos melhores hiperparâmetros, são treinados \textbf{20 modelos independentes}.
    \item O conjunto de teste (20\%) é mantido fixo e só utilizado no final.
    \item Isso permite medir variância e estabilidade por arquitetura.
  \end{itemize}
\end{frame}

\begin{frame}{Métricas e Avaliação}
  \begin{itemize}
  \item Métricas principais:
  \textbf{MCC}, \textbf{F1-score}, \textbf{Accuracy}, \textbf{Specificity}, \textbf{Sensitivity}
  \item Análises:
  \begin{itemize}
    \item Curvas de loss e aprendizado
    \item Boxplots de estabilidade
    \item Interpretação via SHAP e Permutation
  \end{itemize}
\end{itemize}
\end{frame}

\begin{frame}{Explicabilidade e Interpretação}
\begin{itemize}
  \item Técnicas aplicadas:
  \begin{itemize}
    \item \textbf{SHAP}: impacto local de cada feature
    \item \textbf{Permutation Importance}: impacto global
  \end{itemize}

  \item Objetivos:
  \begin{itemize}
    \item Entender decisões do modelo
    \item Destacar features mais relevantes por arquitetura
  \end{itemize}
\end{itemize}
\end{frame}

\begin{frame}{Espaço de Hiperparâmetros}
\textbf{Parâmetros gerais:}
  \begin{itemize}
    \item Dropout: [0.0, 0.5]
    \item Learning rate: [$10^{-4}$, $10^{-2}$] (log)
  \end{itemize}

\vspace{0.8em}
\begin{tabular}{|l|c|l|}
\hline
\textbf{Arquitetura} & \textbf{Camadas} & \textbf{Parâmetros específicos} \\
\hline
\textbf{MLP} & 1–3 densas & 32–512 neurônios por camada \\
\hline
\textbf{CNN1D} & 1–3 convoluções & Kernel: 3/5/7, Filtros: 8–64 \\
\hline
\textbf{LSTM} & 1–2 camadas & Hidden size: 16–256 \\
\hline
\end{tabular}
\end{frame}


\begin{frame}{Pipeline do Experimento}
\begin{enumerate}

\item \textbf{Comandos CLI}
  \begin{itemize}
    \item Escolher Arquitetura (CNN, LSTM, MLP)
  \end{itemize}

  \item \textbf{Divisão dos Dados}
  \begin{itemize}
    \item 80\% treino / 20\% teste
  \end{itemize}

  \item \textbf{Otimização (Optuna)}
  \begin{itemize}
    \item 5-Fold CV e MCC médio
  \end{itemize}

  \item \textbf{Treinamento Final}
  \begin{itemize}
    \item 20 modelos com os melhores hiperparâmetros (por arquitetura)
    \item Avaliação no conjunto de teste
  \end{itemize}

  \item \textbf{Análise}
  \begin{itemize}
    \item Boxplots, curvas, classification report
    \item SHAP e Permutation Importance
  \end{itemize}
\end{enumerate}
\end{frame}

\section{Resultados}

\begin{frame}{Melhores Hiperparâmetros}
\begin{itemize}
  \item \textbf{CNN1D}
  \begin{itemize}
    \item \textbf{Parâmetros:}
      \begin{itemize}
        \item Dropout: 0.1
        \item Learning Rate: 0.00059
        \item Threshold: 0.5
        \item Filtro: 59, Kernel: 6
        \item Camadas Convolucionais: 4
        \item Camadas Densas: 1 com 300
      \end{itemize}
  \end{itemize}
  
  \item \textbf{MLP}
  \begin{itemize}
    \item \textbf{Parâmetros:}
      \begin{itemize}
        \item Dropout: 0.3
        \item Learning Rate: 0.00128
        \item Threshold: 0.6
        \item Camadas Densas: 2 com 643
      \end{itemize}
  \end{itemize}

  \item \textbf{LSTM}
  \begin{itemize}
    \item \textbf{Parâmetros:}
      \begin{itemize}
        \item Dropout: 0.1
        \item Learning Rate: 0.00143
        \item Threshold: 0.7
        \item Camadas LSTM: 1 com 87
      \end{itemize}
  \end{itemize}
\end{itemize}
\end{frame}

\begin{frame}{Importância dos Hiperparâmetros (Optuna)}
  \begin{itemize}
    \item \textbf{Learning rate} é o hiperparâmetro mais relevante em todas as arquiteturas.
    \item \textbf{CNN1D:} kernel size e número de filtros também têm impacto importante.
    \item \textbf{MLP:} número de camadas e taxa de dropout afetam o desempenho.
    \item \textbf{LSTM:} sensível ao número de camadas e unidades escondidas.
  \end{itemize}

  \vspace{0.5em}
  \begin{columns}
    \column{0.33\textwidth}
    \centering
    \includegraphics[width=\linewidth]{data/optuna/param_importance/optuna_param_importance_CNN1D_chest_Sc_4_T_binary_two} \\
    \scriptsize CNN1D
    \column{0.33\textwidth}
    \centering
    \includegraphics[width=\linewidth]{data/optuna/param_importance/optuna_param_importance_MLP_chest_Sc_4_T_binary_two} \\
    \scriptsize MLP
    \column{0.33\textwidth}
    \centering
    \includegraphics[width=\linewidth]{data/optuna/param_importance/optuna_param_importance_LSTM_chest_Sc_4_T_binary_two} \\
    \scriptsize LSTM
  \end{columns}
\end{frame}

\begin{frame}{Convergência do Optuna (Histórico de Trials)}
  \begin{itemize}
    \item Atenção ao dataset!
    \item \textbf{CNN1D} mostrou convergência mais estável e suave.
    \item \textbf{MLP} atingiu ótimos resultados rapidamente com menos oscilações.
    \item \textbf{LSTM} apresentou maior variabilidade entre trials, refletindo instabilidade.
  \end{itemize}

  \vspace{0.5em}
  \begin{columns}
    \column{0.33\textwidth}
    \centering
    \includegraphics[width=\linewidth]{data/optuna/convergencia/optuna_convergencia_CNN1D_chest_Sc_4_T_binary_two} \\
    \scriptsize CNN1D
    \column{0.33\textwidth}
    \centering
    \includegraphics[width=\linewidth]{data/optuna/convergencia/optuna_convergencia_MLP_chest_Sc_4_T_binary_two} \\
    \scriptsize MLP
    \column{0.33\textwidth}
    \centering
    \includegraphics[width=\linewidth]{data/optuna/convergencia/optuna_convergencia_LSTM_chest_Sc_4_T_binary_two} \\
    \scriptsize LSTM
  \end{columns}
\end{frame}

%--- BOXPLOTS: TEXTO
\begin{frame}{Análise da Variabilidade}
  \begin{itemize}
    \item Cada arquitetura foi treinada 20 vezes para avaliar a variabilidade.
    \item A CNN1D obteve o melhor desempenho médio e menor variância.
    \item O MLP mostrou resultados estáveis e competitivos.
    \item O LSTM apresentou maior variabilidade e inconsistência nos resultados.
  \end{itemize}
\end{frame}

\begin{frame}{Boxplot das Métricas – MLP}
  \centering
  \includegraphics[width=0.6\linewidth]{data/boxplots/all/boxplot_MLP_chest_Sc_4_T_binary_two.png}
  
  \vspace{0.5em}
  \textbf{MLP}
\end{frame}

\begin{frame}{Boxplot das Métricas – CNN1D}
  \centering
  \includegraphics[width=0.6\linewidth]{data/boxplots/all/boxplot_CNN1D_chest_Sc_4_T_binary_two.png}
  
  \vspace{0.5em}
  \textbf{CNN1D}
\end{frame}

\begin{frame}{Boxplot das Métricas – LSTM}
  \centering
  \includegraphics[width=0.6\linewidth]{data/boxplots/all/boxplot_LSTM_chest_Sc_4_T_binary_two.png}
  
  \vspace{0.5em}
  \textbf{LSTM}
\end{frame}

\begin{frame}{Boxplot das Métricas}
  \begin{columns}
    \column{0.33\textwidth}
    \centering
    \includegraphics[width=\linewidth]{data/boxplots/all/boxplot_MLP_chest_Sc_4_T_binary_two.png}
    \textbf{MLP}

    \column{0.33\textwidth}
    \centering
    \includegraphics[width=\linewidth]{data/boxplots/all/boxplot_CNN1D_chest_Sc_4_T_binary_two.png}
    \textbf{CNN1D}

    \column{0.33\textwidth}
    \centering
    \includegraphics[width=\linewidth]{data/boxplots/all/boxplot_LSTM_chest_Sc_4_T_binary_two.png}
    \textbf{LSTM}
  \end{columns}
\end{frame}

\begin{frame}{Médias e Desvios Padrão das Métricas}
  \scriptsize
\begin{table}[H]
\centering
\caption{Resultados dos modelos na posição \textbf{chest}, cenário \textbf{Sc\_4\_T}, com rotulagem \textbf{binária (binary\_two)}. Os valores estão no formato \textit{média ± desvio padrão}.}
\begin{tabular}{lcccccc}
\toprule
\textbf{Modelo} & \textbf{MCC} & \textbf{Acurácia} & \textbf{Precisão} & \textbf{Recall} & \textbf{TNR} \\
\midrule
\textbf{CNN1D} & 0.94 ± 0.01 & 0.99 ± 0.00 & 1.00 ± 0.00 & 0.99 ± 0.00 & 0.98 ± 0.02 \\
\textbf{LSTM}  & 0.81 ± 0.06 & 0.96 ± 0.01 & 0.99 ± 0.01 & 0.96 ± 0.01 & 0.91 ± 0.08 \\
\textbf{MLP}   & 0.95 ± 0.02 & 0.99 ± 0.00 & 0.99 ± 0.00 & 0.99 ± 0.01 & 0.96 ± 0.03 \\
\bottomrule
\end{tabular}
\label{tab:resultados_chest_sc4t}
\end{table}
  \begin{itemize}
\normalsize
    \item CNN1D e MLP apresentam alta acurácia e estabilidade.
    \item LSTM demonstra maior variabilidade e desempenho inferior, principalmente em MCC e especificidade.
  \end{itemize}
\end{frame}


\begin{frame}{Curva de Perda — MLP}
  \centering
  \includegraphics[width=0.7\linewidth]{data/loss_curves/MLP_chest_Sc_4_T_binary_two_loss_curve_model_3.png}
  \vspace{0.5em}
\end{frame}

\begin{frame}{Curva de Perda — CNN1D}
  \centering
  \includegraphics[width=0.7\linewidth]{data/loss_curves/CNN1D_chest_Sc_4_T_binary_two_loss_curve_model_14.png}
  \vspace{0.5em}
\end{frame}

\begin{frame}{Curva de Perda — LSTM}
  \centering
  \includegraphics[width=0.7\linewidth]{data/loss_curves/LSTM_chest_Sc_4_T_binary_two_loss_curve_model_16.png}
  \vspace{0.5em}
\end{frame}

\begin{frame}{Curvas de Perda}
  \begin{columns}
    \column{0.33\textwidth}
    \centering
    \includegraphics[width=\linewidth]{data/loss_curves/MLP_chest_Sc_4_T_binary_two_loss_curve_model_3.png}
    \textbf{MLP}

    \column{0.33\textwidth}
    \centering
    \includegraphics[width=\linewidth]{data/loss_curves/CNN1D_chest_Sc_4_T_binary_two_loss_curve_model_14.png}
    \textbf{CNN1D}

    \column{0.33\textwidth}
    \centering
    \includegraphics[width=\linewidth]{data/loss_curves/LSTM_chest_Sc_4_T_binary_two_loss_curve_model_16.png}
    \textbf{LSTM}
  \end{columns}
\end{frame}

\begin{frame}{Comparação - Curva de Loss}

\begin{itemize}
    \item \textbf{CNN1D}
    \begin{itemize}
        \item Treinamento com boa convergência (perda $\approx 0.01$ no final).  
        \item Validação estável até $\sim$20 épocas, mas explosão de perda após $\sim$22 épocas (overfitting evidente).  
    \end{itemize}
    
    \item \textbf{LSTM}  
    \begin{itemize}
        \item Perda de validação consistentemente mais alta que a de treinamento (gap significativo).  
        \item Overfitting desde o início (modelo memoriza rapidamente os dados de treino).  
        \item Possível necessidade de mais dados ou maior regularização.
    \end{itemize}
    
    \item \textbf{MLP}  
    \begin{itemize}
        \item Treinamento estável, mas perda de validação extremamente volátil.  
        \item Picos abruptos (até $\sim10$ de perda) indicam instabilidade do otimizador ou taxa de aprendizado muito alta.  
    \end{itemize}

\end{itemize}

\end{frame}

%--- LEARNING CURVES: TEXTO
\begin{frame}{Curvas de Aprendizado}
  \begin{itemize}
    \item Avalia impacto da quantidade de dados no desempenho.
    \item CNN1D: melhora significativa com mais dados — potencial de escalabilidade.
    \item MLP: robusto mesmo com poucos dados.
    \item LSTM: resultados imprevisíveis, sem tendência clara.
  \end{itemize}
\end{frame}

\begin{frame}{Curva de Aprendizado – MLP}
  \centering
  \includegraphics[width=0.65\linewidth]{data/learning_curves/loss/learning_curve_loss_MLP_chest_Sc_4_T_binary_two.png}
  
  \vspace{0.5em}
  \textbf{MLP}
\end{frame}

\begin{frame}{Curva de Aprendizado – CNN1D}
  \centering
  \includegraphics[width=0.65\linewidth]{data/learning_curves/loss/learning_curve_loss_CNN1D_chest_Sc_4_T_binary_two.png}
  
  \vspace{0.5em}
  \textbf{CNN1D}
\end{frame}

\begin{frame}{Curva de Aprendizado – LSTM}
  \centering
  \includegraphics[width=0.65\linewidth]{data/learning_curves/loss/learning_curve_loss_LSTM_chest_Sc_4_T_binary_two.png}
  
  \vspace{0.5em}
  \textbf{LSTM}
\end{frame}


%--- LEARNING CURVES: IMAGENS
\begin{frame}{Curvas de Aprendizado (Loss)}
  \begin{columns}
    \column{0.33\textwidth}
    \centering
    \includegraphics[width=\linewidth]{data/learning_curves/loss/learning_curve_loss_MLP_chest_Sc_4_T_binary_two.png}
    \textbf{MLP}

    \column{0.33\textwidth}
    \centering
    \includegraphics[width=\linewidth]{data/learning_curves/loss/learning_curve_loss_CNN1D_chest_Sc_4_T_binary_two.png}
    \textbf{CNN1D}

    \column{0.33\textwidth}
    \centering
    \includegraphics[width=\linewidth]{data/learning_curves/loss/learning_curve_loss_LSTM_chest_Sc_4_T_binary_two.png}
    \textbf{LSTM}
  \end{columns}
\end{frame}

\begin{frame}{Learning Curve (Loss)}
\small

\vspace{0.4em}
\begin{itemize}
  \item \textbf{CNN1D:}  
  \begin{itemize}
    \item Perda de validação próxima à de treino na maioria dos pontos.
    \item Pequenas oscilações indicam certa instabilidade, mas não há sinais claros de overfitting.
  \end{itemize}

  \vspace{0.3em}
  \item \textbf{LSTM:}
  \begin{itemize}
    \item Val Loss consistentemente acima do Train Loss.
    \item Indício de overfitting leve a moderado, especialmente nas maiores porcentagens de dados.
  \end{itemize}

  \vspace{0.3em}
  \item \textbf{MLP:}
  \begin{itemize}
    \item Gap considerável entre Train e Val Loss nas menores quantidades de dados.
    \item Com mais dados, as perdas convergem, reduzindo o overfitting.
  \end{itemize}
\end{itemize}

\vspace{0.4em}
\textbf{Conclusão:} Todos os modelos se beneficiam do aumento de dados. LSTM demonstra maior tendência a overfitting. CNN e MLP apresentam boa generalização na maior parte dos cenários.

\end{frame}


\begin{frame}{Curvas de Aprendizado (Métricas)}
  \begin{columns}
    \column{0.33\textwidth}
    \centering
    \includegraphics[width=\linewidth]{data/learning_curves/metrics/learning_curve_metrics_MLP_chest_Sc_4_T_binary_two.png}
    \textbf{MLP}

    \column{0.33\textwidth}
    \centering
    \includegraphics[width=\linewidth]{data/learning_curves/metrics/learning_curve_metrics_CNN1D_chest_Sc_4_T_binary_two.png}
    \textbf{CNN1D}

    \column{0.33\textwidth}
    \centering
    \includegraphics[width=\linewidth]{data/learning_curves/metrics/learning_curve_metrics_LSTM_chest_Sc_4_T_binary_two.png}
    \textbf{LSTM}
  \end{columns}
\end{frame}

\begin{frame}{SHAP}
  \begin{columns}
    \column{0.33\textwidth}
    \centering
    \includegraphics[width=\linewidth]{data/shap/shap_importance_MLP_chest_Sc_4_T_binary_two.png}
    \textbf{MLP}

    \column{0.33\textwidth}
    \centering
    \includegraphics[width=\linewidth]{data/shap/shap_importance_CNN1D_chest_Sc_4_T_binary_two.png}
    \textbf{CNN1D}

    \column{0.33\textwidth}
    \centering
    \includegraphics[width=\linewidth]{data/shap/shap_importance_LSTM_chest_Sc_4_T_binary_two.png}
    \textbf{LSTM}
  \end{columns}
\end{frame}

\begin{frame}{Análise de Importância de Atributos (SHAP)}
\small
\textbf{Objetivo:} Comparar quais sensores (acelerômetro vs giroscópio) foram mais relevantes para cada modelo.

\vspace{0.4em}
\begin{itemize}
  \item \textbf{CNN1D:}
  \begin{itemize}
    \item Acelerômetro dominou a decisão (\texttt{acc\_x} e \texttt{acc\_y} com maior importância).
    \item Giroscópio teve contribuição secundária, relativamente equilibrada entre os eixos.
  \end{itemize}

  \vspace{0.3em}
  \item \textbf{MLP:}
  \begin{itemize}
    \item Forte dependência de um único atributo (\texttt{acc\_z}), seguido de \texttt{gyr\_x}.
    \item Indica sensibilidade a padrões específicos do eixo vertical do acelerômetro.
  \end{itemize}

  \vspace{0.3em}
  \item \textbf{LSTM:}
  \begin{itemize}
    \item Giroscópio foi mais relevante (\texttt{gyr\_y} e \texttt{gyr\_z} com maiores valores SHAP).
    \item Acelerômetro teve menor impacto, sugerindo que o LSTM capturou melhor dinâmicas angulares no tempo.
  \end{itemize}
\end{itemize}

\vspace{0.4em}
\textbf{Conclusão:} Cada arquitetura extraiu padrões distintos: CNN e MLP priorizam aceleração, enquanto LSTM explora mais variações de rotação.

\end{frame}


%--- PERMUTATION IMPORTANCE
\begin{frame}{Permutation Importance}
  \begin{columns}
    \column{0.33\textwidth}
    \centering
    \includegraphics[width=\linewidth]{data/permutation_importance/permutation_importance_CNN1D_chest_Sc_4_T_binary_two}
    \textbf{CNN1D}
    \column{0.33\textwidth}
    \centering
    \includegraphics[width=\linewidth]{data/permutation_importance/permutation_importance_MLP_chest_Sc_4_T_binary_two}
    \textbf{MLP}
    \column{0.33\textwidth}
    \centering
    \includegraphics[width=\linewidth]{data/permutation_importance/permutation_importance_LSTM_chest_Sc_4_T_binary_two}
    \textbf{LSTM}
    
  \end{columns}
\end{frame}


%--- SÍNTESE E CONCLUSÕES

\section{Conclusão e Trabalhos Futuros}

\begin{frame}{Síntese e Conclusões}
  \begin{itemize}
    \item A abordagem proposta mostrou que redes neurais — mesmo arquiteturas simples como \textbf{MLPs} — podem ser eficazes.
    \item O modelo \textbf{CNN1D} e \textbf{MLP} apresentou o melhor desempenho geral:
    \begin{itemize}
      \item Alta acurácia, F1-score, MCC e baixa variância.
      \item Boa interpretabilidade e escalabilidade.
    \end{itemize}
    \item O \textbf{LSTM}, apesar de capturar bem o longo prazo, mostrou instabilidade mesmo após otimização com Optuna. (*)
    \end{itemize}
\end{frame}

%--- LIMITAÇÕES

\begin{frame}{Limitações do Estudo}
  \begin{itemize}
    \item Avaliação restrita ao cenário \texttt{Sc\_4\_T} com sensores na posição \texttt{chest}.
    \item Erro no "shuffle"
    \item Não foi realizado a engenharia de features.
    \item Interpretação dos modelos feita apenas com dados offline (estáticos).
    \item Poucos dados.
  \end{itemize}
\end{frame}

%--- TRABALHOS FUTUROS

\begin{frame}{Trabalhos Futuros}
  \begin{itemize}
    \item Avaliar outros \textbf{dispositivos e posições corporais} (ex: lado esquerdo e direito).
    \item Avaliar ganho com as features em domínio da frequência (Fourier).
    \item Fazer corretamente a engenharia de features para treinar os modelos. (magnitude)
    \item Explicar SHAP temporalmente? ForcePlots, Waterfall, Dependence, etc.
    \item Explorar \textbf{métodos ensemble} e \textbf{conformal prediction} para maior confiabilidade.
    \item Desenvolver uma \textbf{versão embarcável ou em tempo real} para aplicação prática.
    \item Avaliar \textbf{arquiteturas híbridas}.
  \end{itemize}
\end{frame}


\begin{frame}[standout]
  Obrigado! \\
\end{frame}

\end{document}
